
\documentclass[11pt,a4paper]{article}

\usepackage[T1]{fontenc}
\usepackage[utf8]{inputenc}
\usepackage[lithuanian]{babel}
\usepackage{lmodern}
\usepackage{microtype}
\usepackage[a4paper,margin=2.1cm]{geometry}
\usepackage{xcolor}
\usepackage{enumitem}
\usepackage[most]{tcolorbox}
\usepackage{titlesec}
\usepackage{fancyhdr}
\usepackage{parskip}
\usepackage{hyperref}

\definecolor{heading}{HTML}{1E293B}
\definecolor{boxbg}{HTML}{EEF2F7}
\definecolor{boxframe}{HTML}{CBD5E1}

\hypersetup{
  colorlinks=true,
  linkcolor=heading,
  urlcolor=heading
}

\titleformat{\section}
  {\Large\bfseries\color{heading}}
  {}{0pt}{}
\titleformat{\subsection}
  {\large\bfseries\color{heading}}
  {}{0pt}{}

\titlespacing*{\section}{0pt}{1.0em}{0.45em}
\titlespacing*{\subsection}{0pt}{0.8em}{0.35em}

\setlist[itemize]{
  leftmargin=1.5em,
  itemsep=0.15em,
  topsep=0.3em
}

\pagestyle{fancy}
\fancyhf{}
\fancyfoot[C]{\thepage}
\renewcommand{\headrulewidth}{0pt}

\newtcolorbox{infobox}{
  colback=boxbg,
  colframe=boxframe,
  boxrule=0.6pt,
  arc=2mm,
  left=3mm,
  right=3mm,
  top=2.5mm,
  bottom=2.5mm
}

\begin{document}

\begin{center}
{\LARGE\bfseries\color{heading}
Lietuvos komandinė mokinių matematikos olimpiada\\
prof. Jono Kubiliaus taurei laimėti
}

\vspace{0.5em}

{\large\itshape Trumpas gidas pradedančiajam olimpiadininkui}
\end{center}

\vspace{0.8em}

\section{Kas tai per olimpiada?}

\textbf{Lietuvos komandinė mokinių matematikos olimpiada prof. Jono Kubiliaus taurei laimėti}, dažnai trumpai vadinama \textbf{Kubiliaus taure}, – viena svarbiausių komandinių matematikos olimpiadų Lietuvoje.

Ji rengiama Vilniaus universitete ir paprastai vyksta rudenį. Olimpiadą organizuoja Vilniaus universiteto Matematikos ir informatikos fakultetas.

Pagrindinis jos skirtumas nuo Matulionio konkurso ar miesto matematikos olimpiados yra tas, kad čia varžosi ne atskiri mokiniai, o \textbf{komandos}.

Visa komanda gauna vieną bendrą matematinę užduotį ir kartu siekia surinkti kuo daugiau taškų.

\section{Kas sudaro komandą?}

Komandoje gali būti \textbf{iki 5 mokinių}.

Komandos paprastai atstovauja:
\begin{itemize}
  \item gimnazijai;
  \item miestui ar savivaldybei.
\end{itemize}

Tai nėra konkursas, į kurį kiekvienas mokinys tiesiog individualiai užsiregistruoja. Olimpiadoje dalyvauja organizatorių pakviestos komandos.

Kvietimai siejami su ankstesniais matematikos olimpiadų rezultatais, todėl patekimas į Kubiliaus taurės komandą jau savaime reiškia, kad mokinys patenka į gana stiprią olimpiadinę aplinką.

\section{Kiek laiko trunka olimpiada?}

Komanda uždavinius sprendžia \textbf{4 valandas}.

Tačiau šios keturios valandos labai skiriasi nuo keturių valandų individualioje miesto olimpiadoje.

Čia penki mokiniai gali \textbf{dirbti kartu}: tartis, dalintis uždaviniais, tikrinti vieni kitų idėjas ir ieškoti geriausio sprendimo.

\section{Kiek būna uždavinių?}

Kubiliaus taurėje komandai pateikiama net \textbf{20 olimpiadinių uždavinių}.

\begin{infobox}
\centering
\textbf{5 mokiniai \textrightarrow\ 20 uždavinių \textrightarrow\ 4 valandos}
\end{infobox}

Tai labai svarbus skirtumas nuo individualių olimpiadų.

Nesitikima, kad vienas mokinys bandys išspręsti visus 20 uždavinių. Komandos stiprybė yra ta, kad penki žmonės gali dirbti prie skirtingų problemų ir vėliau sujungti savo rezultatus.

\section{Kokie būna uždaviniai?}

Tai tikri \textbf{olimpiadinio pobūdžio uždaviniai}.

Gali pasitaikyti:
\begin{itemize}
  \item algebros;
  \item geometrijos;
  \item skaičių teorijos;
  \item kombinatorikos;
  \item funkcinių lygčių;
  \item nelygybių;
  \item kitų nestandartinio matematinio mąstymo uždavinių.
\end{itemize}

Uždavinių yra daug, todėl jų sunkumas nevienodas.

Vienus stiprus olimpiadininkas gali išspręsti gana greitai, o prie kitų visa komanda gali ilgai ieškoti pagrindinės idėjos.

Svarbu ir tai, kad \textbf{uždavinio numeris nėra patikimas jo sunkumo rodiklis}. Todėl negalima manyti, kad pirmi uždaviniai būtinai bus lengviausi, o paskutiniai – sunkiausi.

\section{Kaip vertinami sprendimai?}

Čia, kaip ir klasikinėje matematikos olimpiadoje, \textbf{vien atsakymo neužtenka}.

Reikia pateikti matematiškai pagrįstą sprendimą.

Už uždavinį paprastai galima gauti:
\begin{itemize}
  \item 5 taškus – už pilną sprendimą;
  \item 4, 3, 2 arba 1 tašką – už įvairaus lygio dalinį sprendimą;
  \item 0 taškų – jei reikšmingos teisingos pažangos nėra.
\end{itemize}

Tačiau Kubiliaus taurė turi vieną labai įdomią vertinimo ypatybę.

Jeigu tam tikrą uždavinį pilnai išsprendžia \textbf{labai mažai komandų}, pilnas jo sprendimas gali būti vertinamas daugiau negu įprastais 5 taškais.

Tai reiškia, kad ypač sunkus ir retai išsprendžiamas uždavinys komandai gali būti vertingesnis.

\section{Kodėl komandinė olimpiada yra kitokia?}

Individualioje olimpiadoje esi vienas:

\begin{center}
\textbf{tavo idėjos \textrightarrow\ tavo sprendimas \textrightarrow\ tavo rezultatas.}
\end{center}

Kubiliaus taurėje atsiranda dar vienas gebėjimas – \textbf{matematinis komandinis darbas}.

Vienas komandos narys gali būti ypač stiprus geometrijoje, kitas – skaičių teorijoje, trečias greitai pastebėti kombinatorines idėjas.

Todėl stipri komanda nebūtinai yra penki vienodi matematikai.

Kartais daug geriau turėti žmones, kurių stipriosios pusės \textbf{papildo viena kitą}.

Taip pat labai svarbu mokėti paaiškinti savo idėją kitam žmogui. Jeigu radai sprendimą, bet negali jo aiškiai perteikti komandos draugui arba užrašyti, komanda iš tavo atradimo gauna daug mažiau naudos.

\section{Kaip nustatomi nugalėtojai?}

Komisija patikrina visus komandų pateiktus sprendimus ir susumuoja gautus taškus.

Komandos rikiuojamos pagal bendrą taškų sumą.

Jeigu kelios komandos surenka tiek pat taškų, pirmiausia žiūrima, kuri komanda turi daugiau \textbf{pilnai išspręstų uždavinių}. Jei ir tada rezultatai lygūs, lyginami kiti sprendimų įvertinimai.

Pirmąją vietą užėmusiai Lietuvos gimnazijos ar savivaldybės komandai atitenka \textbf{profesoriaus Jono Kubiliaus taurė}.

\section{Kuo Kubiliaus taurė ypatinga?}

Kubiliaus taurė nėra tik dar vienas Lietuvos matematikos konkursas.

Ji glaudžiai susijusi su tarptautine komandine matematikos olimpiada \textbf{„Baltijos kelias“}.

Būtent toks pats pagrindinis formatas naudojamas ir „Baltijos kelyje“:

\begin{infobox}
\centering
\textbf{5 mokinių komanda kartu sprendžia 20 uždavinių.}
\end{infobox}

Kubiliaus taurės rezultatai naudojami kviečiant stiprius mokinius į \textbf{Lietuvos komandos atranką į „Baltijos kelio“ olimpiadą}.

Taigi labai geras pasirodymas Kubiliaus taurėje gali tapti vienu iš žingsnių nuo Lietuvos olimpiadinės matematikos iki \textbf{tarptautinių komandinių varžybų}.

\section{Ar Kubiliaus taurė tinka pradedančiajam?}

Tai jau gana aukšto lygio olimpiadinė aplinka.

Skirtingai nuo Matulionio konkurso, čia nėra dalies, kurios pagrindinis tikslas būtų tiesiog greitai rasti trumpus atsakymus.

Reikia mokėti spręsti tikrus olimpiadinius uždavinius ir pagrįsti savo sprendimus.

Tačiau jaunesniam ar dar tik augančiam olimpiadininkui Kubiliaus taurė labai įdomi kaip \textbf{tikslas, į kurį galima augti}.

Kelias galėtų atrodyti maždaug taip:

\begin{infobox}
\centering
\textbf{Matulionio konkursas \textrightarrow\ miesto olimpiada \textrightarrow\ šalies lygio olimpiadinė matematika \textrightarrow\ Kubiliaus taurės komanda \textrightarrow\ „Baltijos kelio“ atranka}
\end{infobox}

Žinoma, realus kiekvieno mokinio kelias gali būti kitoks.

\section{Ką svarbiausia prisiminti pradedančiajam?}

Kubiliaus taurė moko dalyko, kurio individualiose olimpiadose beveik nėra:

\begin{center}
\textbf{matematiką galima kurti kartu.}
\end{center}

Vienas žmogus pastebi idėją.

Kitas suranda jos silpną vietą.

Trečias pasiūlo pataisymą.

Ketvirtas užrašo įrodymą.

Penktas patikrina, ar neliko loginės spragos.

Ir iš penkių atskirų mąstymų atsiranda vienas komandos sprendimas.

Todėl Kubiliaus taurėje svarbu ne tik būti stipriu matematiku. Reikia mokėti \textbf{klausytis, paaiškinti savo mintį, priimti kito žmogaus idėją ir kartu ieškoti sprendimo}.

Būtent tuo ši olimpiada ir išsiskiria iš daugumos kitų matematikos varžybų.

\end{document}
